
\documentclass[12pt,a4paper]{article}

\usepackage[margin=1in]{geometry}
\usepackage{amsmath}
\usepackage{graphicx}
\usepackage{xcolor}
\usepackage{booktabs}
\usepackage{hyperref}

\title{\textbf{Introduction to LaTeX}}
\author{Aakriti}
\date{\today}

\begin{document}

\maketitle

\begin{abstract}
This document demonstrates the basic features of LaTeX,
including formatting, mathematical equations, lists, and tables.
\end{abstract}

\tableofcontents
\newpage

\section{Introduction}

LaTeX is a document preparation system used to create
professional technical documents, reports, and research papers.

\section{Features of LaTeX}

\begin{itemize}
    \item Professional document formatting
    \item Mathematical equations
    \item Tables and figures
    \item Automatic numbering
\end{itemize}

\section{Mathematical Equations}

Pythagoras theorem:

\[
a^2+b^2=c^2
\]

Quadratic formula:

\[
x=\frac{-b\pm\sqrt{b^2-4ac}}{2a}
\]

\section{Ordered List}

\begin{enumerate}
    \item Introduction
    \item Main Content
    \item Conclusion
\end{enumerate}

\section{Table Example}

\begin{table}[h]
\centering
\caption{Student Marks}
\begin{tabular}{ccc}
\toprule
Name & Subject & Marks \\
\midrule
Aman & Mathematics & 90 \\
Riya & Physics & 95 \\
Rahul & Chemistry & 88 \\
\bottomrule
\end{tabular}
\end{table}

\section{Conclusion}

LaTeX helps create structured, professional,
and well-formatted technical documents.

\end{document}
